\begingroup
\raggedbottom
\clubpenalty=10000
\widowpenalty=10000
\chapter{Experiments, Error Analysis, and Scientific Claims}
\label{ch:experiments-claims}

A comparison can answer several different questions. Did one fitted system score higher on these particular cases? Does its fitting procedure perform better across new random starts? Would it improve outcomes for a new population or after a deployment action? Each question concerns a different source of variation and needs different evidence.

Consider a hypothetical course-support assistant that retrieves current syllabus passages and either returns a supported answer or routes the request to a human. A lexical retriever with a template answerer scores $0.842$ on a protected question set; a contextual retriever with a grounded generator scores $0.846$. The difference is $0.004$, or $0.4$ percentage points. These constructed scores describe the observed comparison. They do not determine which component helped, how much a new sample would change the result, or whether the new system handles escalation-required requests well.

\par\addvspace{\intextsep}\noindent\begin{minipage}{\linewidth}
\centering\small
\begin{tabular}{@{}L{5.6cm}L{2.3cm}L{1.8cm}@{}}
\toprule
System & Answer accuracy & Difference \\
\midrule
Lexical retrieval + template & $0.842$ & -- \\
Contextual retrieval + generator & $0.846$ & $+0.004$ \\
\bottomrule
\end{tabular}
\captionof{table}{A constructed pilot comparison. Both systems require the same question set, document snapshot, and grading rule before this difference has a clear meaning.}
\label{tab:score-improved-so-what}
\end{minipage}\par\addvspace{\intextsep}

The same distinction applies to a pedestrian detector. Better daytime average precision is evidence about that evaluated slice, under its localization and matching rules. It does not establish better night-rain performance, lower latency, or safer vehicle control. A useful experiment connects a numerical contrast to a specified claim.

\section{A Comparison Has an Estimand}

\par\noindent\begin{minipage}{\linewidth}
An \emph{estimand} is the quantity the study intends to estimate. For two fixed predictors $f_A,f_B$, a specified case population $P$, and an integrable loss $\ell$, one possible estimand is
\[
\Delta_P=\mathbb E_P[\ell(Y,f_B(X))-\ell(Y,f_A(X))].
\]
\end{minipage}\par\addvspace{\belowdisplayskip}
Positive $\Delta_P$ means A has lower expected loss. For accuracy, use the difference of correctness indicators instead. State the sign convention; a positive loss difference and a positive accuracy difference need not use the same subtraction order.

A comparison of \emph{fitting procedures} has another expectation over their training samples and fitting randomness. A seed-only experiment on one fixed training and test set estimates variation conditional on those sets. A repeated study with newly sampled training sets concerns a broader procedure. Neither expectation can be inferred merely by changing the name of a reported score.

For a fixed test set, its measured difference is a descriptive fact once both predictions and the rubric are fixed. A confidence interval requires a sampling model, such as independent cases drawn from a stated population. Treating a convenience benchmark as a random sample of every future institution silently broadens the claim. A useful local result can remain useful with its scope stated.

\par\addvspace{\intextsep}\noindent\begin{minipage}{\linewidth}
\centering
\begin{tikzpicture}[ml diagram]
\foreach \x/\name/\flowtext in {0/question/question,2.8/protocol/protocol,5.6/contrast/contrast,8.4/claim/claim}{
\node[ml box,minimum width=2.1cm] (\name) at (\x,0) {\flowtext};}
\draw[ml arrow] (question)--(protocol);
\draw[ml arrow] (protocol)--(contrast);
\draw[ml arrow] (contrast)--(claim);
\node[ml label] at (0,-0.65) {population, outcome};
\node[ml label] at (2.8,-0.65) {controls, selection};
\node[ml label] at (5.6,-0.65) {variation, failures};
\node[ml label] at (8.4,-0.65) {supported scope};
\end{tikzpicture}

\captionof{figure}{Define the comparison, protect its observations, estimate the contrast, and state what it supports. Uncertainty and failure analysis inform the contrast and its scope.}
\label{fig:experiment-workflow}
\end{minipage}\par\addvspace{\intextsep}

Before tuning, record the target population, main metric, meaningful baseline, selection process, and consequences that the average metric might omit. This fixes the question before the most flattering result is known. Exploration can revise the question, but a revised claim needs evaluation that accounts for the exploration.

\par\addvspace{\intextsep}\noindent\begin{minipage}{\linewidth}
\centering\small
\begin{tabular}{@{}L{3.0cm}L{6.7cm}@{}}
\toprule
Comparison record & What makes the claim assessable \\
\midrule
Claim and population & Fixed models or fitting procedures; source, time, and eligible cases \\
Outcome and consequences & Primary metric, aggregation, important slices, and severe errors \\
Control and resources & Baseline, information access, fitting and selection budgets \\
Protected protocol & Data versions, grouping, preprocessing, document snapshot, and rubric \\
Variation and inference & Sampling unit, run plan, pairing, interval assumptions, and multiplicity \\
Remaining scope & Untested mechanisms, environments, actions, and alternative explanations \\
\bottomrule
\end{tabular}
\captionof{table}{A comparison record begins before fitting and is updated with the observed contrast, uncertainty, and remaining limitations. It supports an argument rather than a guarantee.}
\label{tab:prerun-claim-sheet}
\label{tab:experiment-claim-sheet}
\end{minipage}\par\addvspace{\intextsep}

For the pilot assistant, the initial claim concerns answer accuracy against a grounded lexical-template baseline on a fixed question population. Faithfulness and escalation remain separate outcomes. Replacing an unsupported answer with safe escalation may help the complete system even when an answer-only rubric penalizes it. The output policy and rubric must agree about what counts as success.

\section{Protect the Comparison and Its Controls}

A reproducible record identifies the data and document versions, eligible cases, grouping and split rules, transformations fitted from training data, candidate settings, stopping rule, and final evaluation. It also records initialization, sampling, decoding, retrieval tie breaking, and any other randomness that can move the result. Exact numerical identity may require additional control of hardware and numerical operations; a broader replication asks whether the substantive finding persists under a documented repeat.

Henderson et al.\ studied variation from seeds and implementation choices in specified deep reinforcement-learning comparisons \citep{henderson2018}. Pineau et al.\ described the NeurIPS reproducibility program and its emphasis on reporting the experimental procedure \citep{pineau2021}. These studies motivate explicit records; they do not establish that every deterministic or stochastic comparison requires the same number of runs.

A shared random seed is a setting, not a proof that two algorithms experience matched randomness. Different architectures may consume random numbers differently. Pair runs through a meaningful common circumstance, such as the same independently chosen data split or prespecified run block, and explain the construction.

\subsection{Baselines Test a Particular Alternative}

A majority-class rule tests whether classification exceeds a prevalence-based shortcut. A linear model tests whether a more flexible feature rule improves the chosen task. Persistence and seasonal-repeat forecasts test whether a temporal model adds value beyond readily available history. A lexical retriever over the same current documents is an informative control for a course-support retrieval claim.

A baseline need not be the hardest conceivable opponent. It must address the alternative explanation at issue. An ungrounded generator can answer a different question from a grounded lexical-template system. Give each candidate the information its comparison permits, and record meaningful differences in labels, pretraining, search effort, inference cost, and latency.

Equal trial counts are not a universal definition of a controlled comparison: one method's trial may cost far more or explore a different search space. A study can compare fixed resource budgets or the best documented procedures affordable in practice. Its claim must include that choice. A heavily tuned new method against a casually configured baseline does not isolate the method from its extra search effort.

Human grading and postprocessing are part of the procedure. Filtering failed outputs from one system alone changes the evaluated population. Evaluating different document snapshots changes the available evidence. Online experimentation studies likewise document how logging and metric-definition problems can distort apparent improvements \citep{kohavi2012,kohavi2013}. Repeating a flawed procedure precisely does not remove the flaw.

\section{Separate Sources of Variation}

Variation across fitting runs, across evaluation cases, and across human judgments has different meanings. More seeds can improve an estimate of conditional average fitting performance. More independent test cases can improve an estimate for new cases. Additional raters can characterize a specified judgment process. Repeating one source of randomness does not automatically reduce the others.

For independent, identically distributed run scores $m_1,\ldots,m_S$ with finite variance and $S>1$,
\[
\bar m=\frac1S\sum_{s=1}^S m_s,\qquad
s_m^2=\frac1{S-1}\sum_{s=1}^S(m_s-\bar m)^2.
\]
The standard deviation describes run variability. The estimated standard error $s_m/\sqrt S$ describes uncertainty in their mean under the stated independence model. It is not the variability of the next individual run. Choosing and reporting only the best run estimates a selected extreme, rather than typical performance.

\subsection{A Constructed Augmentation Tradeoff}

Consider five runs on fixed clean and shifted digit sets. The following accuracies are constructed, rather than measurements establishing robustness to every real camera transformation:
\[
\begin{aligned}
\text{without augmentation, clean: }& .947,.956,.944,.956,.958,\\
\text{without augmentation, shifted: }& .153,.124,.142,.107,.098,\\
\text{with augmentation, clean: }& .871,.882,.833,.882,.878,\\
\text{with augmentation, shifted: }& .844,.858,.856,.873,.844.
\end{aligned}
\]
The means are $0.9522,0.1248,0.8692,0.8550$, respectively; their sample standard deviations are approximately $0.0063,0.0231,0.0207,0.0120$. In these runs, augmentation costs $0.0830$ mean clean accuracy and gains $0.7302$ mean shifted accuracy. Report both changes. Their direction across five runs does not establish performance on untested shift families or represent all possible training samples.

\subsection{A Repeated Pilot on Fixed Questions}

For the support assistant, suppose prespecified paired run blocks give
\[
\begin{aligned}
B:&\quad .842,.839,.847,.841,.844,\\
A:&\quad .846,.845,.850,.844,.848.
\end{aligned}
\]
The means are $0.8426$ and $0.8466$. Paired run differences are
$d=(.004,.006,.003,.003,.004)$, with mean $0.004$ and sample standard deviation about $0.001225$. These runs describe fitting variability on the fixed question set. They do not yet quantify the uncertainty from drawing new student questions or using another policy collection.

A crossed experiment with several fits and several cases produces a table of outcomes, not automatically $S n$ independent observations. Predictions share fitted models, and the same difficult cases recur across fits. An analysis must respect those dependencies. If the intended claim concerns fresh training samples as well as fresh cases, the study needs evidence about both.

\section{Paired Differences and Confidence Intervals}

When two fixed systems are evaluated on the same cases, let $d_i$ be A's correctness minus B's correctness. Each paired observation belongs to one case. Six cases with differences $(0,1,-1,0,0,1)$ give $\bar d=1/6$: two wins, one loss, and three ties. The loss still matters if it concerns an escalation-required request.

Pairing retains where predictions disagree. Its variance is
\[
\operatorname{Var}(A_i-B_i)=\operatorname{Var}(A_i)+\operatorname{Var}(B_i)
-2\operatorname{Cov}(A_i,B_i).
\]
Positive covariance from shared case difficulty often improves precision relative to independently sampled scores. Zero or negative covariance can eliminate or reverse that advantage. Pairing is appropriate because of the design; it is not an unconditional variance-reduction theorem.

For independent paired differences with sample mean $\bar d$ and sample standard deviation $s_d$, an exact interval under a normal population model is
\[
\bar d\pm t_{1-\alpha/2,n-1}\frac{s_d}{\sqrt n}.
\]
For a sufficiently well-behaved large sample with finite variance, this can be an approximation \citep{nistPairedCI}. A few differences in $\{-1,0,1\}$ are not a normal sample merely because the formula is available. Subject, source, or temporal dependence also changes the appropriate sampling unit.

For the five paired pilot runs, using $t_{.975,4}\approx2.776$ gives the illustrative interval
\[
.004\pm2.776\frac{.001225}{\sqrt5}
\approx[.00248,.00552].
\]
It relies on a plausible normal model for independent run-block differences and concerns their conditional mean. Five runs give little information about that model's tails. This interval should not be presented as an interval for arbitrary new questions.

Now suppose two \emph{fixed} fitted systems score $0.842$ and $0.846$ on $1{,}000$ independent questions: A alone is correct on $28$, B alone on $24$, and they tie on $948$. The paired mean is still $0.004$, but $s_d\approx0.2281$, giving estimated standard error $0.00721$. A normal approximation gives about $[-0.0101,0.0181]$. This constructed case comparison has a different uncertainty from the seed comparison, despite the identical mean difference.

For binary correctness and an independent paired-case null of equal win and loss probabilities, condition on the number of discordant cases. The number of A-only wins then has a binomial distribution with success probability $1/2$. This yields an exact paired test under those assumptions. It tests a null about disagreements, not safety or practical importance, and it does not make a selected benchmark representative.

A frequentist $95\%$ confidence procedure covers the fixed population estimand in $95\%$ of repeated samples under its assumptions. The observed interval is not, without a separate probabilistic model, a $95\%$ posterior probability that the parameter lies inside it. A $p$-value likewise concerns outcomes at least as incompatible with a specified null as the observed statistic; it is not the probability that the null is true.

\subsection{Bootstrap the Independent Unit}

A paired nonparametric bootstrap samples $n$ observed pairs with replacement and recalculates the contrast. For a mean difference, this is equivalent to sampling the $d_i$ values. The percentile interval takes the empirical $\alpha/2$ and $1-\alpha/2$ quantiles of many resampled means \citep{efron1993bootstrap}.

Its coverage is approximate and needs suitable sampling and regularity conditions. A tiny, skewed sample or an inadequately observed tail can defeat the approximation. More bootstrap replicates reduce simulation error, not uncertainty caused by an inadequate original sample. Resample subjects or suitable blocks when that is the independent unit; separately resampling two systems destroys their pairing. Not every positively or negatively dependent dataset is repaired by arbitrary blocks.

For nonlinear metrics such as F1 or average precision, recompute the entire metric on each resampled paired dataset. Averaging invented per-case F1 contributions would estimate a different quantity. The bootstrap also retains benchmark selection and labeling problems. It estimates variation under a model; it does not validate the target population.

\section{Cross-Validation Evaluates a Procedure}

In $K$-fold cross-validation, each fold is scored by a predictor fitted using the other folds. All learned preprocessing and selection must remain on the training side of the relevant evaluation boundary. Groups or chronology replace ordinary random folds when the intended use requires them. For unequal fold sizes, a pooled case-average loss weights each fold by its size; an unweighted mean of fold metrics can answer a different question.

Hyperparameter search belongs inside the outer training portion. Nested cross-validation uses an inner split or loop to choose settings and an outer held-out fold to evaluate the resulting selection procedure. It estimates performance for that documented procedure at the outer training size. It neither proves that the chosen setting is the best possible configuration nor gives an exact guarantee for a final model refitted on all the data.

If a team chooses the best of twelve configurations by their common outer-fold scores and reports that maximum as an untouched evaluation, those folds have become development evidence. A new protected evaluation or a selection-aware analysis is needed for the new claim. Nested selection prevents this particular reuse when correctly implemented; it does not cure biased sources, mislabeled targets, or cross-group leakage.

Fold scores share overlapping training sets. They are not independent replicates, and their sample standard deviation divided by $\sqrt K$ is not automatically a valid standard error. Bengio and Grandvalet show that there is no universal unbiased variance estimator for $K$-fold cross-validation without further assumptions \citep{bengio2004cv}. Repeating folds on the same observations adds partition variation, not new independent people. Report that spread for what it measures.

\section{Power, Search, and Multiple Comparisons}

An interval can exclude zero while the gain is too small to justify added latency or cost. A substantial useful gain can remain uncertain when few independent units are available. Before running the study, specify the smallest effect that would matter and compare it with the protocol's ability to detect that effect.

For two independent conditions with $n$ independent scores each and a common known standard deviation $\sigma$, the standard error of their mean difference is $\sigma\sqrt{2/n}$. A normal planning approximation for a two-sided level-$\alpha$ test with power $1-\beta$ is
\[
\Delta_{\rm MDE}\approx(z_{1-\alpha/2}+z_{1-\beta})\sigma\sqrt{2/n},
\qquad
n\approx\frac{2(z_{1-\alpha/2}+z_{1-\beta})^2\sigma^2}{\Delta^2}.
\]
The approximation uses the alternative's likely tail and is most useful away from very low power; an exact normal calculation includes both rejection tails. Estimated variances, small-sample $t$ tests, multiplicity, and dependence require further adjustment.

With $\alpha=.05$, power $.8$, $z_{.975}+z_{.8}\approx2.80$, $\sigma=.01$, and target difference $.004$, the planning calculation gives about $98$ independent runs \emph{per condition}. It concerns conditional seed variation if those are the independent units. For $n$ paired differences, replace $\sigma\sqrt{2/n}$ by $\sigma_d/\sqrt n$; a meaningful pairing and its covariance change the required budget. An effect below the planned minimum detectable effect is harder to detect reliably, not necessarily zero or attributable to luck.

Search creates additional opportunities to select a favorable outcome. If $m$ independent valid tests all have true nulls and rejection probability $\alpha$, the probability of at least one false rejection is
\[
1-(1-\alpha)^m.
\]
For $m=20$ and $\alpha=.05$, this is about $0.642$. The approximation $m\alpha$ is useful only when the combined probability remains small. Independence is required for the product formula; under arbitrary dependence the union bound gives at most $m\alpha$, capped at one.

Bonferroni tests each of $m$ prespecified hypotheses at level $\alpha/m$, controlling the probability of any false rejection through that union bound when each null test is valid. False-discovery-rate procedures address a different quantity, the expected fraction of false rejections among rejections, with zero when nothing is rejected. Their assumptions must also be stated. A correction to a displayed list of tests does not account automatically for an undisclosed adaptive search or invalid individual tests.

Prespecify primary comparisons and report the search that produced the final system. Development exploration is useful, but trying prompts or thresholds against the final test and hiding all but the winner changes the evidence. More computation cannot restore a held-out boundary after its results influenced design.

\section{Ablation Measures a Specified Change}

An ablation removes or changes a component while keeping the remaining comparison as controlled as the question requires. Removing a module from a fixed fitted system measures reliance on that module at inference. Refitting a system without it measures the performance of a different fitting procedure. Retuning the ablated procedure changes the estimand again. State which intervention was performed.

Suppose a constructed follow-up benchmark compares two retrievers and two answerers under the same task definition:

\par\addvspace{\intextsep}\noindent\begin{minipage}{\linewidth}
\centering\small
\begin{tabular}{@{}L{5.0cm}L{2.5cm}L{2.2cm}@{}}
\toprule
System & Answer accuracy & Faithfulness pass \\
\midrule
Lexical retrieval + template & $.842$ & $.989$ \\
Contextual retrieval + template & $.861$ & $.990$ \\
Lexical retrieval + generator & $.848$ & $.944$ \\
Contextual retrieval + generator & $.872$ & $.962$ \\
\bottomrule
\end{tabular}
\captionof{table}{A constructed factorial comparison on a follow-up benchmark, distinct from the pilot. The outcome definition and denominators must be consistent across systems.}
\label{tab:support-assistant-ablation}
\end{minipage}\par\addvspace{\intextsep}

Switching retrieval changes accuracy by $.019$ with the template and $.024$ with the generator. Switching answerer changes it by $.006$ with lexical retrieval and $.011$ with contextual retrieval. The difference of differences is $.005$: these observed changes are not additive. The full system gains $.030$ over the lexical-template system while losing $.027$ faithfulness pass rate. Uncertainty and the severity of unsupported claims are still needed to assess that tradeoff.

The controlled changes support comparisons of these implemented components and their interaction under this protocol. They do not uniquely identify every internal mechanism or establish that a generator generally helps. If several settings, costs, or data sources change with the component, those differences belong to the intervention being evaluated.

\section{Error Analysis Generates Testable Explanations}

Suppose routine questions have $95$ correct answers out of $100$, while escalation-required questions have $15$ successful responses out of $30$ under an explicit rubric. If these slices are disjoint and exhaust the sample, overall success is $110/130\approx.846$, while the second slice is only $.5$. The aggregate weights the sample's prevalence, not the consequences of each failure. Overlapping slices require separate counts; adding their denominators can count a case twice.

Reading errors can reveal unsupported additions, stale passages, missed escalation cues, and false-confidence patterns. Those observations suggest explanations. They do not establish that a representation lacks understanding or that one suspected cue caused the failure. A useful next experiment compares alternatives that would produce different observable outcomes.

\par\addvspace{\intextsep}\noindent\begin{minipage}{\linewidth}
\centering\small
\begin{tabular}{@{}L{3.0cm}L{3.2cm}L{3.5cm}@{}}
\toprule
Observed pattern & Explanation to test & Informative comparison \\
\midrule
Wrong passages on paraphrases & Lexical mismatch or incomplete index & Matched retrieval methods with index coverage recorded \\
Unsupported answer conditions & Answer policy or missing evidence & Answerers given the same retrieved passages \\
Weak multilingual slice & Coverage, labeling, or measurement differences & Protected slice with annotation and source controls \\
Night-rain detector misses & Visibility, source shift, or cue dependence & Matched capture tests and specified interventions \\
\bottomrule
\end{tabular}
\captionof{table}{A failure pattern motivates competing explanations and a comparison that can narrow them. It does not identify a unique mechanism by itself.}
\label{tab:error-cluster-next-experiment}
\end{minipage}\par\addvspace{\intextsep}

Slices discovered by inspecting final-test errors become development evidence if they guide a fix. Preserve the distinction between describing those failures and claiming the revised system generalizes better. A fresh protected sample can test the revised hypothesis.

\subsection{Training Gaps Are Diagnostic Clues}

The squared-loss bias--variance decomposition in Section~\ref{sec:bias-variance} averages over specified training randomness at a specified input distribution. One training error and one test error do not reveal its terms. Poor optimization, different populations, mislabeled outcomes, regularization, and sampling variation can produce similar patterns.

\par\addvspace{\intextsep}\noindent\begin{minipage}{\linewidth}
\centering\small
\begin{tabular}{@{}L{1.5cm}L{1.5cm}L{3.0cm}L{3.7cm}@{}}
\toprule
Train loss & Test loss & Hypotheses & Evidence to distinguish them \\
\midrule
High & High & Restricted fit, optimization, or noise & Fit diagnostics, target quality, and controlled capacity changes \\
Low & High & Overfitting, shift, or leakage differences & Repeated fits, matched populations, and split review \\
High & Low & Unequal difficulty or evaluation behavior & Compare losses, populations, and training-only transformations \\
Low & Low & Good observed fit & Quantify uncertainty and evaluate the intended conditions \\
\bottomrule
\end{tabular}
\captionof{table}{Training and test losses suggest investigations. They are not measurements of bias or variance and need comparable definitions.}
\label{tab:bias-variance-diagnostic}
\end{minipage}\par\addvspace{\intextsep}

More suitable data, different regularization, or richer features can help under some explanations. Their effects depend on the fitting procedure and the new data. Additional observations from the same biased source need not fix a missing deployment population. A smaller train--test gap alone is not the optimization target; a model can reduce the gap by getting worse on both sets.

\section{Human Evaluation Defines Another Measurement}

A reviewer needs a specified question, an eligible population, and enough evidence to judge it. Randomize presentation order and conceal system identity where possible. A reviewer who sees only the answer can judge fluency while missing an unsupported policy condition; a reviewer who sees the source can assess support. The information provided to the reviewer changes the outcome measured.

\par\addvspace{\intextsep}\noindent\begin{minipage}{\linewidth}
\centering\small
\begin{tabular}{@{}L{2.3cm}L{3.7cm}L{3.7cm}@{}}
\toprule
Dimension & Rubric question & Example failure \\
\midrule
Faithfulness & Are substantive claims supported by permitted evidence? & Added eligibility condition \\
Escalation & Was an authority-requiring request routed correctly? & Unsupported direct advice \\
Completeness & Are the necessary next step and boundary included? & Missing office or deadline \\
Clarity & Can the reader act without guessing? & Fluent but ambiguous action \\
\bottomrule
\end{tabular}
\captionof{table}{Separate judgments support separate claims. A preference score does not automatically measure factual support or safe routing.}
\label{tab:human-eval-rubric}
\end{minipage}\par\addvspace{\intextsep}

Record disagreement and adjudication instead of assuming that a majority vote establishes truth. Raters can share systematic errors or a mistaken rubric. Ratings from one person across many cases, or many ratings of one case, also share dependencies. The uncertainty analysis should match the intended population of cases and reviewers. A language variety can affect rubric interpretation as well as model behavior.

\section{Predictive Evidence and Intervention Effects}

A predictor estimates an association, perhaps a conditional mean or distribution. Its accuracy does not establish what would happen after changing an input or an institutional policy. Students who seek tutoring may differ in prior difficulty, schedules, or motivation. Neither a positive nor a negative association between visits and grades determines the effect of offering tutoring.

For a binary intervention $T$, write potential outcomes $Y(1),Y(0)$. The average treatment effect is
\[
\tau=\mathbb E[Y(1)-Y(0)].
\]
Under consistency, randomized assignment independent of both potential outcomes, positive assignment probabilities, and an appropriately defined intervention without interference between units,
\[
\mathbb E[Y\mid T=1]-\mathbb E[Y\mid T=0]=\tau.
\]
The equality follows because each group reveals its assigned potential outcome and randomization gives the same potential-outcome distribution in both groups. Observational comparisons need additional identification assumptions; a predictive fit does not supply them \citep{pearl2009}. Randomization of an offer estimates the effect of that offer, not automatically the effect of attendance among those who choose to attend.

A detector can predict pedestrian presence well while the complete braking policy fails because of latency, trajectory uncertainty, or fallback behavior. Evaluating the predictor and evaluating a closed-loop action are different studies. An offline score does not establish the consequences of deploying the action policy.

A report can therefore make a bounded claim: the selected assistant achieved the measured improvement on the protected sample under the documented procedure, with specified uncertainty and slice results. Broader statements about new institutions, safe automation, or the internal mechanism require their own evidence. A negative result, a visible tradeoff, or a narrower operating scope can be the substantive finding.

\begin{studyquestions}
\item How do a fixed-model case comparison and a comparison across fitting runs differ?
\item What must a baseline control, and why are equal trial counts not enough to define every resource comparison?
\item Why can shared cases improve paired precision, and when might they fail to do so?
\item What randomness does a confidence interval describe, and what does its coverage mean?
\item Why are cross-validation folds and repeated ratings not automatically independent units?
\item What do an ablation and a difference of differences establish under a specified protocol?
\item Why can a smaller training--test gap or a larger aggregate score conceal worse behavior?
\item Which assumptions turn a randomized intervention comparison into evidence for a causal effect?
\end{studyquestions}

\begin{exercises}[State the estimand, unit, and sign convention.]
\item \exercisetag{Design}{Intermediate}Write a local model-comparison claim and a stronger claim the same study would not support. Name the population, protocol, and uncertainty needed for the local claim.
\item \exercisetag{Calculation}{Core}For run accuracies $.81,.84,.79,.82$, compute the sample mean, sample standard deviation, and estimated standard error under independent identically distributed runs. Distinguish the last two quantities.
\item \exercisetag{Concept}{Core}Give a meaningful classification baseline and a forecasting baseline. State which alternative explanation each addresses.
\item \exercisetag{Design}{Intermediate}Construct a disjoint two-slice example with a strong overall success rate and an unacceptable rare-slice rate. Show the aggregation and define unacceptable behavior.
\item \exercisetag{Design}{Intermediate}Plan an ablation of pretraining, augmentation, and postprocessing. Distinguish removal from a fixed fit, refitting, and retuning, and include a possible interaction.
\item \exercisetag{Concept}{Intermediate}State a broader claim still unsupported after beating a baseline on one benchmark.
\item \exercisetag{Concept}{Intermediate}For a support assistant, contrast a slice requiring supported answer generation with one requiring correct escalation. Specify the evidence and rubric for each.
\item \exercisetag{Diagnosis}{Advanced}A detector improves familiar daytime AP but worsens night-rain and unfamiliar-camera AP. State the supported comparison and what further evidence a deployment decision needs.
\item \exercisetag{Design}{Intermediate}Complete the comparison record in Table~\ref{tab:experiment-claim-sheet} for one predictor. Identify which fields can be fixed before results and which require observations.
\item \exercisetag{Calculation}{Advanced}For paired differences $(-1,0,0,1,1)$, compute their mean and sample standard deviation. Give one size-five bootstrap resample and its mean. Explain which assumptions fail if the cases form one correlated sequence.
\item \exercisetag{Calculation}{Intermediate}Two paired scores each have variance $.01$ and covariance $.008$. Compute the difference variance. Repeat with covariance $-.008$, and compare with independently sampled scores.
\item \exercisetag{Calculation}{Intermediate}Reproduce the five-run pilot's paired mean, sample standard deviation, and illustrative $95\%$ interval using $t_{.975,4}=2.776$. State which uncertainty this interval omits.
\item \exercisetag{Calculation}{Intermediate}For twenty independent true-null tests at level $.05$, calculate the chance of at least one rejection. Give the Bonferroni threshold for familywise level $.05$ and explain what dependence changes.
\item \exercisetag{Calculation}{Intermediate}Compute both retrieval effects, both generator effects, and the accuracy interaction in Table~\ref{tab:support-assistant-ablation}. Compute the full system's faithfulness change from the lexical-template system.
\item \exercisetag{Calculation}{Advanced}Using the normal planning approximation, compute runs per condition for $\sigma=.01,\Delta=.004$ and quantile sum $2.80$. If paired differences have standard deviation $.005$, compute the paired-run requirement and state the additional pairing assumptions.
\end{exercises}

\clearpage
\begin{challengeexercises}[Make selection and identification assumptions explicit.]
\item \exercisetag{Diagnosis}{Advanced}A team's best run beats a baseline by $0.4$ percentage points. Give three surviving explanations and a study design that would narrow them without reusing the final test.
\item \exercisetag{Diagnosis}{Advanced}A model improves overall performance but worsens a critical subgroup. Explain how to report the result and which evidence should govern a deployment recommendation.
\item \exercisetag{Design}{Advanced}Identify the hardest field of your comparison record to defend. Design the additional evidence needed and state what would remain unestablished.
\item \exercisetag{Concept}{Advanced}A result survives strong controls and repeated runs but has no ablation or slice evaluation. Distinguish its remaining performance claim from unsupported mechanism and deployment claims.
\item \exercisetag{Diagnosis}{Advanced}A team repeats five-fold cross-validation ten times on the same people and divides the standard deviation of fifty fold scores by $\sqrt{50}$. Explain why this is not automatically a standard error, and design an evaluation respecting subjects and selection.
\item \exercisetag{Proof}{Advanced}Derive the randomized-treatment equality from potential outcomes. Explain how noncompliance, interference, or observational assignment changes the estimand or invalidates the stated identification.
\end{challengeexercises}

\clearpage
\begingroup\pagestyle{empty}\cleardoublepage\endgroup
\endgroup
